\section{Kimi Linear y KDA: actualización de la memoria con granularidad más fina}

Esta sección entra en KDA (Kimi Delta Attention), el módulo central de Kimi Linear. En el programa se menciona que KDA parte de trabajos anteriores como Gated DeltaNet y sustituye la decay de granularidad gruesa por una decay de granularidad más fina: antes, las distintas dimensiones de un mismo attention head compartían una sola tasa de decaimiento; ahora cada dimensión puede tener su propia tasa de decaimiento. La intuición es que el hidden state, que es limitado, exprese mejor memorias de distintas escalas temporales.

\begin{figure}[H]
\centering
\includegraphics[width=0.96\textwidth]{kimi-linear-kda.png}
\caption{Kimi Linear / KDA: Kimi Delta Attention aprovecha mejor la memoria con una decay de granularidad más fina. Diagrama conceptual propio, elaborado a partir de la conversación de 00:14:39--00:18:56.}
\end{figure}

\begin{knowledgebox}{Lectura de la figura: KDA es una nueva combinación de ideas anteriores}
Gated DeltaNet aporta la base, la decay de granularidad gruesa asegura la eficiencia, la decay de granularidad fina aumenta la capacidad expresiva, KDA reorganiza la actualización de la memoria y, al final, todo se integra en una arquitectura hybrid. Muchas arquitecturas nuevas no se inventan de la nada: recombinan mecanismos anteriores en el contexto de hardware nuevo y escalas nuevas.
\end{knowledgebox}

\subsection{Scaling Ladder: superar la escala pequeña antes de ampliarla}

Esta subsección explica cómo Kimi reduce el riesgo de prueba y error en la arquitectura. La atención lineal no se lleva directamente a un modelo grande en cuanto se escribe la fórmula; tiene que validarse escala por escala.

Internamente, Kimi utiliza una Scaling Ladder: un módulo se compara primero con Full Attention a escala pequeña, y solo si su desempeño es suficientemente bueno pasa a una escala mayor para seguir haciendo scale. El proceso se parece a superar los niveles de un juego y evita gastar cómputo en prueba y error directamente a la escala máxima. El diseño de KDA también se convirtió poco a poco en candidato después de varias rondas de selección entre formas de combinación y de ampliación de escala.

\begin{figure}[H]
\centering
\includegraphics[width=0.96\textwidth]{scaling-ladder.png}
\caption{Scaling Ladder: tras superar la escala pequeña se pasa a una escala mayor, con comparación frente a Full Attention. Diagrama conceptual propio, elaborado a partir de la conversación de 00:18:56--00:20:20.}
\end{figure}

\begin{knowledgebox}{Lectura de la figura: la arquitectura no se decide de improviso}
La selección a escala pequeña reduce costos; cuando las métricas alcanzan el nivel requerido, se pasa a una escala mayor y se vuelve a comparar con la línea base de Full Attention. Solo si el costo, el desempeño y el comportamiento de scaling resultan satisfactorios puede volverse candidato para producción.
\end{knowledgebox}

\subsection{La proporción de mezcla tres a uno}

A continuación se trata el problema de la proporción en la arquitectura hybrid. Decidir cuántas capas globales conservar y cuántas sustituir por capas lineales consiste, en esencia, en encontrar entre la capacidad expresiva y la eficiencia un punto de ingeniería que pueda escalarse.

En la entrevista se menciona que Kimi Linear inserta una capa de atención completa por cada tres capas de KDA, y que la proporción tres a uno se ha ido volviendo un consenso. En esencia, esta proporción usa Full Attention para mantener la capacidad expresiva global y KDA para reducir el costo de inferencia en la mayoría de las capas. No es un teorema matemático, sino un compromiso empírico surgido de los experimentos de ingeniería.

\begin{knowledgebox}{Asimilación de los mecanismos de KDA}
\begin{tabularx}{\textwidth}{p{0.22\textwidth}p{0.34\textwidth}X}
\toprule
Mecanismo & Función & Por qué importa \\
\midrule
Delta Rule & Actualiza el estado de memoria de forma recursiva & Da a la attention una capacidad de actualización de estado similar a la de una RNN. \\
Gating & Controla qué información se conserva y cuál se olvida & Permite que el modelo elija qué información histórica debe permanecer. \\
Decay & Tasa de decaimiento de la memoria & Distintas informaciones pueden tener distintas escalas temporales. \\
Fine-grained decay & Cada dimensión tiene su propio decaimiento & Aprovecha mejor un hidden state limitado. \\
Capas globales hybrid & Insertan Full Attention de forma periódica & Garantizan como respaldo la expresión global y la interacción de largo alcance. \\
\bottomrule
\end{tabularx}
\end{knowledgebox}

\begin{warningbox}{KDA no se reduce a «quitar la Softmax»}
Si solo se quita la Softmax, el modelo puede quedarse sin suficiente capacidad expresiva. El valor de KDA está en rediseñar la actualización del estado, el gating y el decaimiento para que las capas lineales conserven en lo posible la información histórica útil; las capas de Full Attention se encargan de completar la interacción global.
\end{warningbox}

\subsection{Resumen de la sección}

Lo esencial de Kimi Linear no es «sustituir todo por atención lineal», sino usar KDA dentro de una arquitectura hybrid para actualizar la memoria de forma eficiente y garantizar un nivel mínimo de desempeño con unas pocas capas globales. La Scaling Ladder es el mecanismo que permite llevar este diseño, paso a paso, hasta el entrenamiento a gran escala.
